(1) 令刚体获得正向冲量 \(J\)，则粒子碰后 \(v'=v-J/m\)，刚体质心速度 \(V=J/M\)。接触点力矩为 \(-bJ\)，角速度变为 \(\omega=-bJ/I\)；接触点沿冲量方向速度为 \(V-b\omega=J(1/M+b^2/I)\)。

(2) Newton恢复条件为 \(J(1/m+1/M+b^2/I)-v=ev\)，从而 \(J=(1+e)v/[m^{-1}+M^{-1}+b^2/I]\)。再代入可求 \(v'、V、\omega\)；当 \(b\ne0\) 时即使冲量水平也必激发刚体自旋。

(3) 线动量变化为 \(-J+MV=0\)。绕碰撞瞬间原刚体质心的角动量变化，粒子产生 \(+bJ\)，刚体产生 \(I\omega=-bJ\)，故两者相消。偏心项 \(b^2/I\) 增加相对接触点的逆质量，导致相同恢复条件下冲量减小。

(4) 用 \(K'=mv'^2/2+MV^2/2+I\omega^2/2\)，整理为 \(\Delta K=-m_{\rm eff}(1-e^2)v^2/2\)，其中 \(m_{\rm eff}^{-1}=1/m+1/M+b^2/I\)。\(e=1\) 时动能不损失，\(e=0\) 时碰后接触点法向相对速度为零，\(b=0\) 时恢复一维两质体碰撞。
